\documentclass[a4paper,11pt]{article}
\usepackage[utf8]{inputenc}
\usepackage[spanish]{babel}
\usepackage{amsmath,amssymb}
\usepackage{geometry}
\usepackage{hyperref}
\usepackage{listings}
\usepackage{xcolor}
\usepackage{tcolorbox}
% Generador de apuntes: integración con la biblioteca oficial WEC.
% La paleta y las macros visuales (titulo, secciongrande, seccion,
% subtituloA/B, resumen, lista, pasos, definicion, nota, tabla...) proceden
% de bibliotecas_latex/cabecera.tex. Aquí se sincronizan los colores
% del resaltado de código y el índice automático clicable con hyperref.
\providecommand{\cmd}[1]{\texttt{\detokenize{#1}}}
\providecommand{\phantomsection}{}
\providecommand{\subnota}[2]{\nota{#1}{#2}}
\providecommand{\subnotasin}[2]{\nota{#1}{#2}}
\AtBeginDocument{\ifcsname spanishdeactivate\endcsname\spanishdeactivate{<>}\fi}
\providecommand{\subt}[1]{\par\Needspace{4\baselineskip}\vspace{2mm}\noindent{\bfseries\color{primary}#1}\par\vspace{1mm}}
\providecolor{codeForeground}{HTML}{243247}
\providecolor{codeString}{HTML}{8A4B08}
\providecolor{codeComment}{HTML}{58745C}
\providecolor{codeBackground}{HTML}{F1F7FF}
\providecolor{codeKeyword}{HTML}{2457A6}
\providecolor{consoleBackground}{HTML}{F1F4F8}
\providecolor{primary}{HTML}{0B4F6C}
\providecolor{coloractual}{HTML}{0B4F6C}
\providecommand{\subtituloreal}{primary}
\providecommand{\subtitulocolor}{primary}
\providecommand{\pickcolor}[1]{}

\newcommand{\actualizarColoresCodigo}{%
  \colorlet{codeBackground}{\coloractual!7!white}%
  \colorlet{codeKeyword}{\coloractual!80!black}%
}
\newcommand{\iniciarPaletaSecciones}{\pickcolor{\coloractual}}
\AddToHook{cmd/titulo/after}{\iniciarPaletaSecciones\actualizarColoresCodigo}
\AddToHook{cmd/secciongrande/after}{\actualizarColoresCodigo}

% --- CORRECCIÓN DE TÍTULOS MULTILÍNEA EN BARRAS OFICIALES WEC ---
\ifcsname sectionheader\endcsname
  \ifcsname bartitleboxlen\endcsname\else\newlength{\bartitleboxlen}\fi
  \renewcommand{\sectionheader}[4][primary]{%
    \edef\sectionheadercolor{#1}%
    \setlength{\barleftlen}{#2}%
    \settowidth{\bartextlen}{\headingfontxb#4\color{\sectionheadercolor} #3}%
    \setlength{\barrightlen}{\dimexpr\linewidth - \barleftlen - \bartextlen - 0.9cm\relax}%
    \ifdim\barrightlen>0.7cm
      \noindent\barrule[\sectionheadercolor]{\barleftlen}\hspace{4mm}{\headingfontxb#4\color{\sectionheadercolor} #3}\hspace{4mm}\barrule[\sectionheadercolor]{\barrightlen}\par
    \else
      \setlength{\bartitleboxlen}{\dimexpr\linewidth - 2\barleftlen - 8mm\relax}%
      \noindent\parbox[c]{\barleftlen}{\barrule[\sectionheadercolor]{\barleftlen}}\hspace{4mm}%
      \parbox[c]{\bartitleboxlen}{\centering\headingfontxb#4\color{\sectionheadercolor} #3}\hspace{4mm}%
      \parbox[c]{\barleftlen}{\barrule[\sectionheadercolor]{\barleftlen}}\par
    \fi
  }
\fi

% --- ÍNDICE AUTOMÁTICO CLICABLE CON HYPERREF (PRESERVANDO ESTILOS OFICIALES WEC) ---
\ifcsname secciongrande\endcsname
  \AddToHook{cmd/secciongrande/before}{\phantomsection}
\else
  \DeclareDocumentCommand{\secciongrande}{o m}{%
    \par\Needspace{8\baselineskip}%
    \phantomsection
    \addcontentsline{toc}{section}{#2}%
    \vspace{4mm}\noindent{\fontsize{15}{18}\selectfont\bfseries\color{primary}#2}\par
    \nopagebreak\vspace{2mm}\nopagebreak
  }
\fi

\ifcsname seccion\endcsname
  \AddToHook{cmd/seccion/before}{\phantomsection}
\else
  \DeclareDocumentCommand{\seccion}{o m}{%
    \par\Needspace{6\baselineskip}%
    \phantomsection
    \addcontentsline{toc}{section}{#2}%
    \vspace{3mm}\noindent{\fontsize{13}{16}\selectfont\bfseries\color{primary}#2}\par
    \nopagebreak\vspace{2mm}\nopagebreak
  }
\fi

\ifcsname subtituloA\endcsname
  \let\origsubtituloA\subtituloA
  \renewcommand{\subtituloA}[1]{%
    \phantomsection
    \addcontentsline{toc}{subsection}{#1}%
    \origsubtituloA{#1}%
  }
\else
  \DeclareDocumentCommand{\subtituloA}{m}{%
    \par\Needspace{5\baselineskip}%
    \phantomsection
    \addcontentsline{toc}{subsection}{#1}%
    \vspace{3mm}\noindent{\fontsize{12}{15}\selectfont\bfseries\color{primary}#1}\par
    \nopagebreak\vspace{2mm}\nopagebreak
  }
\fi

\ifcsname subtituloB\endcsname
  \let\origsubtituloB\subtituloB
  \renewcommand{\subtituloB}[1]{%
    \phantomsection
    \addcontentsline{toc}{subsubsection}{#1}%
    \origsubtituloB{#1}%
  }
\else
  \DeclareDocumentCommand{\subtituloB}{m}{%
    \par\Needspace{4\baselineskip}%
    \phantomsection
    \addcontentsline{toc}{subsubsection}{#1}%
    \vspace{2mm}\noindent{\fontsize{11}{14}\selectfont\bfseries\color{primary}#1}\par
    \nopagebreak\vspace{2mm}\nopagebreak
  }
\fi

\lstdefinelanguage{JavaScript}{
  sensitive=true,
  morekeywords={const,let,var,of,async,await,undefined,null,true,false,function,return,if,else,for,while,switch,case,break,continue,new,class,extends,import,export,default,try,catch,throw},
  morecomment=[l]{//},
  morecomment=[s]{/*}{*/},
  morestring=[b]",
  morestring=[b]'
}
\ifcsname codigo\endcsname\else
\newtcblisting{codigo}[1][Java]{listing only,breakable,colback=codeBackground,colframe=\coloractual!28!white,boxrule=0pt,arc=3pt,left=6pt,right=6pt,top=4pt,bottom=4pt,listing options={language=#1,basicstyle=\ttfamily\footnotesize\color{codeForeground},keywordstyle=\color{codeKeyword}\bfseries,morekeywords={const,let,var,of,async,await,undefined,null,true,false,console,document,window,localStorage,fetch,Promise,script,html,head,body,noscript,div,span,p,ul,li,input,form,button,type,src,id,class},stringstyle=\color{codeString},commentstyle=\color{codeComment}\itshape,showstringspaces=false,breaklines=true,columns=fullflexible,keepspaces=true,tabsize=4,upquote=true}}
\fi
\ifcsname consola\endcsname\else
\newtcblisting{consola}{listing only,breakable,colback=consoleBackground,colframe=primary!22!white,boxrule=0pt,arc=3pt,left=6pt,right=6pt,top=4pt,bottom=4pt,listing options={basicstyle=\ttfamily\footnotesize\color{codeForeground},breaklines=true,columns=fullflexible,keepspaces=true,upquote=true}}
\fi

% --- Fallbacks para macros oficiales de la biblioteca en entornos mínimos ---
\providecommand{\paletaAzul}{}
\providecommand{\usarfondo}{}
\providecommand{\titulo}[2]{\begin{center}{\LARGE\bfseries #1}\\[3mm]{\large\bfseries #2}\end{center}\vspace{4mm}}
\providecommand{\coloractual}{blue}
\providecommand{\primary}{blue}
\providecommand{\headingfont}{\bfseries}
\providecommand{\headingfontxb}{\bfseries}
\providecommand{\resumen}[1]{\par\vspace{2mm}\noindent\fbox{\parbox{\dimexpr\linewidth-2\fboxsep-2\fboxrule}{\textbf{Resumen:} #1}}\par\vspace{2mm}}
\providecommand{\definicion}[2]{\par\vspace{2mm}\noindent\fbox{\parbox{\dimexpr\linewidth-2\fboxsep-2\fboxrule}{\textbf{#1:}\ #2}}\par\vspace{2mm}}
\providecommand{\definicionpar}[2]{\par\vspace{2mm}\noindent\fbox{\parbox{\dimexpr\linewidth-2\fboxsep-2\fboxrule}{\textbf{#1:}\ #2}}\par\vspace{2mm}}
\providecommand{\nota}[2]{\par\vspace{2mm}\noindent\fbox{\parbox{\dimexpr\linewidth-2\fboxsep-2\fboxrule}{\textbf{#1:}\ #2}}\par\vspace{2mm}}
\providecommand{\alertas}[1]{\par\vspace{2mm}\noindent\fbox{\parbox{\dimexpr\linewidth-2\fboxsep-2\fboxrule}{\textbf{Alerta:}\ #1}}\par\vspace{2mm}}
\providecommand{\rectangulos}[1]{#1}
\ifcsname lista\endcsname\else
\ExplSyntaxOn
\DeclareDocumentCommand{\lista}{m}{%
  \begin{itemize}%
    \seq_set_split:Nnn \l_tmpa_seq { | } { #1 }%
    \seq_map_inline:Nn \l_tmpa_seq {%
      \tl_set:Nn \l_tmpa_tl { ##1 }%
      \tl_trim_spaces:N \l_tmpa_tl
      \tl_if_empty:NF \l_tmpa_tl {%
        \item \l_tmpa_tl
      }%
    }%
  \end{itemize}%
}
\ExplSyntaxOff
\fi
\providecommand{\pasos}[1]{\begin{enumerate}#1\end{enumerate}}
\providecommand{\tabla}[2][]{#2}
\providecommand{\inputcarpeta}[1]{%
  \immediate\write18{find "#1" -name "*.tex" 2>/dev/null | sort -V | sed 's/.*/\\input{&}/' > "/tmp/in-#1.tex"}%
  \IfFileExists{/tmp/in-#1.tex}{\input{/tmp/in-#1.tex}}{}%
}


\geometry{top=2.5cm, bottom=2.5cm, left=2.5cm, right=2.5cm}

\title{\textbf{ojo}\\ \large Tema 1: Selección de arquitecturas y herramientas de programación}
\author{Generador Automático de Apuntes}
\date{\today}

\begin{document}
\maketitle
\tableofcontents
\newpage

\secciongrande{Modelos de ejecución de código: Los cuatro campos del diseño web}

\resumen{
  El desarrollo web moderno se fundamenta en el modelo cliente-servidor y en cuatro campos interdependientes del diseño: el visual, el tecnológico, el de distribución y el de propósito. Desde los orígenes del HTML estático hasta las aplicaciones dinámicas actuales, el programador debe dominar un campo multidisciplinar que abarca desde la captación de necesidades hasta la presentación final, garantizando estándares mantenidos por el W3C.
}

\subtituloA{Origen y alcance de la web}

En sus orígenes, la creación de una página web requería únicamente el lenguaje HTML, cuyas marcas permitían construir documentos sencillos que cualquier navegador podía visualizar. Estas páginas se alojaban en servidores conectados permanentemente, transmitiéndose a los equipos clientes que las solicitaban.

La web se fundamenta en el concepto de hipertexto: palabras o frases que, al ser seleccionadas, conducen a información relacionada. Ante el auge de esta tecnología, se creó el consorcio W3C, organización encargada de desarrollar y mantener los estándares web que garantizan la interoperabilidad entre diferentes plataformas y navegadores.

Un programador web debe dominar todas las fases de un proyecto: captar las necesidades del cliente, definir objetivos y estrategia, desarrollar la interfaz gráfica, elaborar prototipos y realizar la presentación final. Es imprescindible conocer el orden secuencial de cada fase, ya que el desarrollo web constituye un campo inherentemente multidisciplinar.

\subtituloA{Los cuatro campos del diseño web}

El diseño de un sitio web se articula en torno a cuatro campos fundamentales que determinan su éxito y funcionalidad:

\definicion{Campo visual}{Aspecto estético y de comunicación de la página web. Su objetivo principal es atraer a los clientes y causar una buena sensación en los visitantes, aumentando así las probabilidades de que vuelvan a visitar el sitio. Engloba la usabilidad, la accesibilidad y la experiencia de usuario.}

\definicion{Campo tecnológico}{Decisión sobre cómo se implementará el sitio y qué tecnologías se utilizarán. De esta elección dependen las características fundamentales del sitio: si será una página estática o dinámica, qué sistemas de bases de datos se emplearán, qué lenguajes de servidor se utilizarán y qué frameworks se aplicarán.}

\definicion{Campo de distribución}{Aspecto relativo a la velocidad y fiabilidad de la transmisión de datos entre el servidor y los clientes. Incluye la optimización del ancho de banda, el uso de redes de distribución de contenido y la garantía de disponibilidad del servicio durante las 24 horas del día.}

\definicion{Campo de propósito}{Objetivo fundamental que persigue el sitio web. En la mayoría de los casos la finalidad es económica, buscando alcanzar una gran difusión del producto en el mercado laboral, aunque también puede ser informativa, como la difusión de actos o eventos de carácter general que puedan interesar al público.}

\subtituloA{El modelo cliente-servidor}

Hoy en día, la arquitectura de los sitios web se basa en el modelo cliente-servidor, que consta de tres elementos comunes esenciales:

\lista{
  \textbf{Servidor:} Equipo donde se instalan los elementos hardware y software, las técnicas incrustadas y los lenguajes de programación. Para implementar estas tecnologías web se aconseja utilizar un lenguaje de programación que soporte multihilos, permitiendo atender múltiples solicitudes simultáneas de manera eficiente. |
  \textbf{Máquina cliente:} Equipo del usuario donde se ejecutan los programas que realizan la petición de solicitud web al servidor a través de los navegadores. Estos navegadores deben soportar las tecnologías del sitio web: HTML, CSS, JavaScript y controles ActiveX. |
  \textbf{Conectividad de red:} Infraestructura y características utilizadas para la transmisión de datos entre cliente y servidor, incluyendo protocolos de comunicación, seguridad y velocidad de conexión.
}

El objetivo fundamental de un navegador es presentar la información, escrita normalmente en lenguaje HTML y alojada en un servidor con el servicio web instalado, interpretando y renderizando el contenido para el usuario final.

\nota{Para el examen}{
¿Dónde tendremos implementados los diferentes programas para realizar peticiones de solicitud en el servidor?

\textbf{Respuesta:} En la máquina cliente, concretamente en el navegador web, que es el programa encargado de realizar las solicitudes al servidor. Además, ¿el funcionamiento de un navegador se encarga de alojar una página web para responder a las distintas solicitudes recibidas?

\textbf{Respuesta:} Falso. El navegador no aloja la página; es el servidor quien la aloja y sirve, mientras que el navegador solicita, descarga y presenta dicha información al usuario.
}

\secciongrande{Lenguajes de programación en entorno cliente dos entornos de ejecución}

\resumen{
  La programación web requiere diferenciar entre entornos cliente y servidor. El entorno cliente se ejecuta en el navegador del usuario, mientras que el servidor procesa peticiones y gestiona bases de datos. La elección del lenguaje depende de las ventajas e inconvenientes, siendo crucial la complementariedad entre ambos entornos para sitios dinámicos.
}

\definicion{Entorno cliente}{Entorno de ejecución donde el código es procesado por el dispositivo del usuario, típicamente un navegador web.}

\definicion{Entorno servidor}{Entorno de ejecución donde el código reside y se procesa en el servidor remoto, gestionando lógica de negocio y bases de datos.}

\subtituloA{Diferenciación y elección de lenguajes}

\lista{
  \textbf{Entorno cliente:} Ejecutado por el cliente en el navegador web, enfocado en la interactividad del usuario. |
  \textbf{Entorno servidor:} Ejecutado en el servidor, responsable del almacenamiento de datos y la lógica de negocio compleja.
}

\nota{Para el examen}{
Es vital que el desarrollador evalúe las ventajas e inconvenientes de cada lenguaje, ya que la decisión final es contextual y depende de su pericia profesional.
}

\subtituloA{Complementariedad cliente y servidor}

|
  \textbf{Interacción: |
  Los lenguajes de ambos entornos se complementan en sitios dinámicos, no entorpecen.
  \textbf{Flujo de trabajo:
} El cliente realiza peticiones; el servidor las recibe, procesa (ej. validación de formularios) y almacena datos en bases de datos.
}

\subtituloA{Tecnologías recomendadas en el cliente}

|
  \textbf{HTML y CSS: |
  Lenguajes fundamentales para dar estructura y formato/estilo a la página.
  \textbf{JavaScript:
} Lenguaje clave para crear páginas dinámicas y gestionar la interactividad en el lado del cliente, soportado por la mayoría de navegadores actuales.
}

\subtituloA{Capas de desarrollo web en el cliente con JavaScript}

, el material bruto de la página. \textbf{Capa de estructura:
} Definida por el documento HTML y el modelo de objeto del documento (DOM). \textbf{Capa de presentación:} Definida por CSS, controla el estilo visual.

\textbf{Capa de comportamiento:} Gestionada por JavaScript, controla la interactividad y la lógica del cliente.

\nota{Para el examen}{
El orden jerárquico de las capas en el cliente es: contenido \rightarrow estructura \rightarrow presentación \rightarrow comportamiento (de abajo hacia arriba).
}

\secciongrande{Mecanismos de ejecución de código en un navegador web la arquitectura del navegador}

\resumen{
  La interacción en un navegador comienza con la solicitud de un recurso y finaliza con su visualización. El proceso se adapta a prioridades como velocidad, fidelidad al contenido o seguridad. La mayoría de los navegadores comparten una arquitectura de referencia compuesta por componentes esenciales.
}

\definicion{Arquitectura de referencia}{Conjunto de componentes básicos comunes a la mayoría de los navegadores web, independientemente de su implementación específica.}
\subtituloA{Ciclo de interacción del usuario}

\lista{
  \textbf{Inicio:} El usuario indica la dirección del recurso (url o uri). |
  \textbf{Finalización:} Visualización exitosa del contenido solicitado. |
  \textbf{Dependencias:} El proceso se configura según el propósito (ej. rapidez, fidelidad, seguridad de la comunicación con el servidor).
}

\lista{
  \textbf{Interfaz de usuario:} Gestiona barras de herramientas, progreso de carga, preferencias y se comunica con el sistema operativo. |
  \textbf{Motor del buscador:} Carga urls, soporta navegación básica (anterior/siguiente) y gestiona alertas de javascript. |
  \textbf{Motor de renderizado:} Presenta visualmente el recurso, interpretando html, xml, css y contenido embebido. |
  \textbf{Subsistema de comunicaciones:} Implementa protocolos de transferencia de archivos (ej. http). |
  \textbf{Intérprete de javascript:} Analiza y ejecuta código javascript, incluyendo eventos de ratón o teclado. |
  \textbf{Parser xml:} Carga la representación en árbol (dom) del documento html en memoria para facilitar el acceso a sus elementos. |
  \textbf{Componente de visualización:} Ofrece funcionalidades específicas para la presentación del contenido html. |
  \textbf{Subsistema de persistencia de datos:} Almacena historial, sesiones de usuario, preferencias, certificados de seguridad y cookies.
}
\nota{Para el examen}{
Diferencie entre el motor del buscador y el motor de renderizado. El primero gestiona la carga de la dirección, mientras que el segundo se encarga de la presentación visual del contenido interpretado.
}

\subtituloA{Motores específicos}

\lista{
  \textbf{V8:} Utilizado por Chrome, Edge, Opera y Brave; también soporta el entorno de servidor node.js. |
  \textbf{SpiderMonkey:} Motor de Firefox, desarrollado inicialmente por Brendan Eich en 1995. |
  \textbf{JavaScriptCore:} Motor exclusivo de Safari, compatible con dispositivos ios (iphone/ipad).
}
\subtituloB{Optimización de ejecución}

|
  \textbf{Compilación JIT: |
  Los motores modernos compilan a código máquina justo antes de la ejecución, logrando alta velocidad. \textbf{Actualización:
} Microsoft abandonó Chakra en 2020; Edge utiliza V8 actualmente.

}
\nota{Para el examen}{
¿Por qué el javascript moderno es rápido? Porque los motores como V8 y SpiderMonkey emplean compilación JIT, no son meros intérpretes lentos. Recuerde que Edge ya no usa Chakra.}

\secciongrande{Capacidades y limitaciones de ejecución: ¿por qué JavaScript en el cliente}

\resumen{
  Aunque la gestión dinámica principal ocurre en el servidor, JavaScript es esencial en el cliente para mejorar la usabilidad. Sus capacidades incluyen la preparación de datos, la navegación sin dependencia total de JS, y la actualización dinámica del contenido recibido del servidor.
}

\subtituloA{Razones para usar JavaScript en el cliente}

\lista{
  \textbf{Preparación de datos:} JavaScript no puede escribir archivos en el servidor; su función es preparar datos para su envío o cederlos a otro lenguaje del servidor. |
  \textbf{Navegación resiliente:} El sitio debe ser navegable sin JavaScript; este se usa para modernizar la experiencia, pero no debe ser un requisito absoluto para la funcionalidad básica. |
  \textbf{Actualización dinámica:} Permite examinar datos recibidos del servidor y actualizar selectivamente secciones de la página web sin recargarla completamente.
}

\subtituloA{Soluciones específicas que ofrece JavaScript}

|
  \textbf{Feedback al usuario: |
  Proporcionar respuesta inmediata desde formularios o hipertexto. \textbf{Interfaz amigable:
} Ofrecer una experiencia de usuario moderna y fluida. \textbf{Control de elementos:} Gestionar plugins, marcos de navegación o applets (ej.

Java). \textbf{Validación previa:} Procesar datos de formularios antes de enviarlos al servidor. \textbf{Manipulación del DOM:} Modificar dinámicamente estilos y contenidos en el navegador.}

\subtituloA{Limitaciones de ejecución: El entorno sandbox}

|
  \textbf{Restricción de archivos: |
  JavaScript no puede leer ni escribir archivos locales del disco duro del usuario. \textbf{Seguridad del navegador:
} Todo JS se ejecuta en una caja de arena (sandbox), impidiendo lanzar programas o tocar la configuración del navegador.

\textbf{Política del mismo origen (SOP):} Un script de un dominio no puede acceder a recursos de otro dominio (ej. com). El origen se define por protocolo + dominio + puerto. \textbf{CORS:} El error CORS al acceder a APIs ajenas no es un fallo del código cliente, sino la aplicación de la política SOP; el servidor debe autorizar mediante cabeceras.

\textbf{Vulnerabilidad XSS:} El riesgo existe si se inserta texto no limpiado de un usuario (ej. %%VERBATIM\_BLOCK\_0%%), permitiendo la ejecución de scripts maliciosos. La protección recae en la sanitización de datos por parte del desarrollador.

}
\nota{Para el examen}{
Es falso que HTML gestione la dinámica completa; la dinámica requiere lenguajes de script en el servidor y JavaScript en el cliente para la usabilidad. La protección del navegador no cubre la mala calidad de los datos introducidos por el usuario (XSS).}

\secciongrande{Programación de scripts: Dónde y cómo colocar los scripts}

\resumen{
  Los scripts de javascript se insertan en un documento html utilizando la etiqueta \texttt{<script>}. Pueden ubicarse tanto en la sección \texttt{<head>} como en el cuerpo \texttt{<body>}. La etiqueta de apertura debe incluir el atributo \texttt{type="text/javascript"} para indicar al navegador el formato del código.
}

\definicion{Etiqueta script}{Etiqueta html contenedora que designa un bloque de código javascript, comenzando con \texttt{<script>} y terminando con \texttt{</script>}.}

\subtituloA{Ubicación en el documento html}

\lista{
  \textbf{Encabezado del script:} Colocación entre \texttt{<head>} y \texttt{</head>}. |
  \textbf{Cuerpo del script:} Colocación entre \texttt{<body>} y \texttt{</body>}.
}

\subtituloA{Sintaxis y atributos}

\definicion{Atributo type}{Atributo que especifica el formato del script; en este caso, \texttt{type="text/javascript"}, indicando que el contenido es texto organizado como javascript.}

\begin{codigo}
<script type="text/javascript">
// Código javascript aquí
</script>
\end{codigo}



\nota{Para el examen}{
¿Cuál es la función del atributo \texttt{type="text/javascript"} en la etiqueta \texttt{<script>}?

\textbf{Respuesta:} Comunica al navegador que el contenido es código javascript y no html.
}

\subtituloA{Ejemplo de ejecución}

\begin{codigo}
document.write("Hola Onliner!");
\end{codigo}



\subtituloA{Terminación del script}

\definicion{Terminación de script}{La etiqueta \texttt{</script>} finaliza el bloque de javascript, indicando al navegador que debe volver a procesar el código html subsiguiente.}

\begin{codigo}
<script>
document.write("Hola Onliner!");
</script>
\end{codigo}



\nota{Para el examen}{
¿Qué función se utiliza para escribir texto directamente en la página web desde javascript?

\textbf{Respuesta:} \texttt{document.write()}. ¿Qué símbolo marca el final de una instrucción en javascript?

\textbf{Respuesta:} El punto y coma (\texttt{;}).
}

\secciongrande{Integración del código con las etiquetas html: Integración online o en archivo externo}

\subtituloA{Métodos de integración de javascript}

\definicion{Integración online}{Consiste en incluir el código javascript directamente dentro de las etiquetas %%VERBATIM\_BLOCK\_0%% utilizando el atributo type="texto/javascript"}
\definicion{Integración en archivo externo}{Consiste en referenciar un archivo javascript externo (con extensión .js) mediante el atributo src de la etiqueta %%VERBATIM\_BLOCK\_1%%}
\lista{
  \textbf{Integración online:} %%VERBATIM\_BLOCK\_2%% para evitar problemas de análisis sintáctico.
}
\definicion{Licencias de código}{Permiten añadir licencias, como Creative Commons, para indicar cómo otros desarrolladores pueden utilizar el código.}
\nota{Para el examen}{
¿Cuál es la forma más recomendable de integrar código javascript y por qué?

\textbf{Respuesta:} El archivo externo, debido a la modularidad, facilidad de depuración y rapidez de carga.
}

\secciongrande{Herramientas de programación sobre clientes web: Elección del editor de código}

\resumen{
  La elección de una herramienta de desarrollo depende del programador, pudiendo ser un editor de texto simple. Es crucial seleccionar desarrolladores de código y no editores de contenido como FrontPage. Las pautas incluyen soporte para codificación de colores, verificación de sintaxis y opciones como clientes FTP. Editores gratuitos recomendados son Notepad++, Eclipse, NetBeans, KompoZer y Bluefish.
}

\definicion{Editor de código}{Herramienta especializada en la escritura, edición y gestión de código fuente, optimizada para lenguajes de programación, a diferencia de los editores de contenido orientados a la presentación.}
\subtituloA{Pautas para la selección de herramientas}

\lista{
  \textbf{No utilizar editores de contenido:} Evitar herramientas como FrontPage o Dreamweaver, enfocándose en desarrolladores de código. |
  \textbf{Facilidad para desarrollo web:} Buscar codificación de colores, verificación de sintaxis, inserción de comentarios, generación de blogs y clientes FTP.
}
\nota{Para el examen}{
Diferenciar entre un editor de código (enfocado en la sintaxis) y un editor de contenido (enfocado en la presentación visual).
}

\definicion{Entorno cliente}{El entorno de ejecución del código web (HTML, CSS, JavaScript) en el navegador del usuario, independiente de la conexión a internet para la codificación inicial.}
\subtituloA{Pasos recomendados para un portal web}

\lista{
  \textbf{Codificación:} Escribir código HTML, JavaScript y CSS en el documento. |
  \textbf{Ejecución:} Guardar con extensiones correctas y ejecutar en navegadores; recargar para ver cambios. |
  \textbf{Validación:} Paso crucial para asegurar la validez del código y la consistencia entre navegadores, siguiendo los estándares W3C.
}
\nota{Para el examen}{
La estandarización del consorcio W3C rige las validaciones web; el validador se encuentra en http://validator.w3.org/
}

\definicion{Minificación}{Proceso de eliminar espacios y acortar nombres en el código para reducir el tamaño del archivo y mejorar la velocidad de carga, no para ocultar la lógica.}
\subtituloA{Principio de seguridad web}

\lista{
  \textbf{Exposición del código:} El navegador descarga el código completo, haciéndolo visible a cualquiera. |
  \textbf{Validación de seguridad:} Debe realizarse obligatoriamente en el servidor; la validación en el cliente es solo para comodidad del usuario. |
  \textbf{Claves sensibles:} Nunca deben incluirse en el código del cliente (ej. claves de API, contraseñas).
}
\nota{Para el examen}{
¿Por qué una clave de API no debe estar en el código cliente? Porque el navegador descarga el código completo, exponiendo la clave públicamente.
}

\definicion{Modelo cliente-servidor}{Arquitectura web donde el cliente (navegador) solicita recursos y el servidor los aloja y responde, complementando lenguajes de cliente (HTML, CSS, JS) y de servidor.}
\subtituloA{Componentes del navegador}

\lista{
  \textbf{Arquitectura de referencia:} El navegador incluye ocho componentes, como interfaz, buscador, renderizado, comunicaciones, intérprete JS, parser XML, visualización y persistencia. |
  \textbf{Integración de scripts:} Se usa la etiqueta <script> con rutas relativas, absolutas o remotas, complementada con <noscript> y sintaxis CDATA en XHTML.
}
\nota{Para el examen}{
Recordar que JavaScript es la capa de comportamiento de la web y que el sitio debe funcionar sin depender de él. El flujo es: editor $\rightarrow$ prueba en navegadores $\rightarrow$ validación W3C.
}


\end{document}